\section{Matemática}

\subsection{MDC / GCD}

Algoritmo de Euclides.

\textbf{STL:} \texttt{std::gcd(a, b)}

\textbf{Complexidade:} $O(\log(\min(a,b)))$.

\begin{lstlisting}
long long gcd(long long a, long long b) {
    while (b) {
        a %= b;
        swap(a, b);
    }

    return a;
}
\end{lstlisting}


\subsection{MMC / LCM}

\[
\operatorname{lcm}(a,b)
=
\frac{a}{\gcd(a,b)} \cdot b
\]

\textbf{STL:} \texttt{std::lcm(a, b)}

\textbf{Complexidade:} $O(\log(\min(a,b)))$.

\begin{lstlisting}
long long lcm(long long a, long long b) {
    return a / gcd(a, b) * b;
}
\end{lstlisting}


\subsection{Combinatória}

Combinações distintas escolhendo $k$ elementos dentre $n$ elementos, sem importar a ordem:
\[
\binom{n}{k}
=
\frac{n!}{k!(n-k)!}
\]

Propriedade útil:
\[
\binom{n}{k}
=
\binom{n}{n-k}
\]

Para valores pequenos, pode ser calculado iterativamente evitando
calcular os fatoriais completos.

\begin{lstlisting}
long long C(int n, int k) {
    if (k > n - k)
        k = n - k;

    long long res = 1;

    for (int i = 1; i <= k; i++)
        res = res * (n - k + i) / i;

    return res;
}
\end{lstlisting}


\subsection{Permutação}

Número de maneiras de escolher e ordenar $k$ elementos dentre $n$:
\[
P(n,k)
=
\frac{n!}{(n-k)!}
=
n(n-1)\cdots(n-k+1)
\]

\textbf{Complexidade:} $O(k)$.

\begin{lstlisting}
long long P(int n, int k) {
    long long res = 1;

    for (int i = 0; i < k; i++)
        res *= n - i;

    return res;
}
\end{lstlisting}


\subsection{Fast Pow}

Calcula $a^b$ usando exponenciação binária.

\textbf{Complexidade:} $O(\log b)$.

\begin{lstlisting}
long long binpow(long long a, long long b, long long mod) {
    // Se o problema nao requer mod, basta remover o parametro e o mod do codigo
    a %= mod;
    long long res = 1;
    
    while (b) {
        if (b & 1) res = res * a % mod;
        a = a * a % mod;
        b >>= 1;
    }
    return res;
}
\end{lstlisting}


\subsection{Inverso Modular}

$ x $ é o inverso modular de $a$ módulo $m$ se:
\[
a x \equiv 1 \pmod m
\]

Existe se e somente se:
\[
\gcd(a,m)=1
\]

Se $m$ é primo:
\[
a^{-1} \equiv a^{m-2} \pmod m
\]

\textbf{Complexidade:} $O(\log m)$.

\begin{lstlisting}
long long modInverse(long long a, long long mod) {
    return binpow(a, mod - 2, mod);
}
\end{lstlisting}

\textbf{Requer:} $m$ primo e $a \not\equiv 0 \pmod m$.


\subsection{Totient de Euler}

$\varphi(n)$ = quantidade de inteiros em $[1,n]$
coprimos com $n$.

Se:
\[
n = p_1^{a_1}p_2^{a_2}\cdots p_k^{a_k}
\]

então:
\[
\varphi(n)
=
n
\prod_{p \mid n}
\left(1-\frac{1}{p}\right)
\]

\textbf{Complexidade:} $O(\sqrt n)$.

\begin{lstlisting}
long long phi(long long n) {
    long long res = n;

    for (long long p = 2; p * p <= n; p++) {
        if (n % p == 0) {
            while (n % p == 0)
                n /= p;

            res -= res / p;
        }
    }

    if (n > 1)
        res -= res / n;

    return res;
}
\end{lstlisting}

\subsection{Primos}

\subsubsection{Crivo de Eratóstenes}

Encontra todos os primos até $N$.

\textbf{Complexidade:} $O(N \log\log N)$.

\begin{lstlisting}
vector<bool> isPrime(N + 1, true);
isPrime[0] = isPrime[1] = false;

for (int i = 2; i * i <= N; i++) {
    if (!isPrime[i])
        continue;

    // j = i * i eh uma otimizacao, 
    // tenha em mente que em alguns casos pode ser necessario iniciar com i + i
    for (int j = i * i; j <= N; j += i)
        isPrime[j] = false;
}
\end{lstlisting}


\subsubsection{Fatoração Prima}

Decompõe $n$ em seus fatores primos.

\textbf{Complexidade:} $O(\sqrt n)$.

\begin{lstlisting}
vector<pair<long long, int>> factorize(long long n) {
    vector<pair<long long, int>> factors;

    for (long long p = 2; p * p <= n; p++) {
        if (n % p != 0)
            continue;

        int cnt = 0;

        while (n % p == 0) {
            n /= p;
            cnt++;
        }

        factors.push_back({p, cnt});
    }

    if (n > 1)
        factors.push_back({n, 1});

    return factors;
}
\end{lstlisting}


\subsubsection{Contagem de Divisores}

Se:
\[
n = p_1^{a_1}p_2^{a_2}\cdots p_k^{a_k}
\]

então a quantidade de divisores positivos de $n$ é:
\[
d(n) = (a_1+1)(a_2+1)\cdots(a_k+1)
\]

\textbf{Complexidade:} $O(\sqrt n)$.

\begin{lstlisting}
long long divisorCount(long long n) {
    long long res = 1;

    for (long long p = 2; p * p <= n; p++) {
        if (n % p != 0)
            continue;

        int cnt = 0;

        while (n % p == 0) {
            n /= p;
            cnt++;
        }

        res *= cnt + 1;
    }

    if (n > 1)
        res *= 2;

    return res;
}
\end{lstlisting}